\documentclass[aspectratio=169,10pt]{beamer}
\usetheme[numbering=fraction,progressbar=frametitle]{metropolis}
\usepackage{fontspec}
\setsansfont{Latin Modern Sans}
\usepackage{amsmath,amssymb,mathtools,graphicx,booktabs}
\usepackage{polyglossia}
\setdefaultlanguage{french}
\setbeamertemplate{navigation symbols}{}
\setbeamerfont{frametitle}{size=\large}
\setbeamerfont{normal text}{size=\normalsize}
\setlength{\parskip}{.45em}
\hypersetup{pdfauthor={},pdfsubject={Lecture slides},pdfkeywords={},pdftitle={Dérivation, convexité et continuité}}
\providecommand{\N}{\mathbb{N}}
\providecommand{\Z}{\mathbb{Z}}
\providecommand{\D}{\mathbb{D}}
\providecommand{\Q}{\mathbb{Q}}
\providecommand{\R}{\mathbb{R}}
\providecommand{\C}{\mathbb{C}}
\providecommand{\vect}[1]{\overrightarrow{#1}}
\providecommand{\vi}{\vec{\imath}}
\providecommand{\vj}{\vec{\jmath}}
\providecommand{\vk}{\vec{k}}
\providecommand{\norm}[1]{\left\lVert #1\right\rVert}
\providecommand{\vc}[2]{\begin{pmatrix}#1\\#2\end{pmatrix}}
\providecommand{\deter}[2]{\det\left(#1,#2\right)}
\providecommand{\abs}[1]{\left\lvert #1\right\rvert}
\providecommand{\dd}{\,\mathrm{d}}
\providecommand{\e}{\mathrm{e}}
\providecommand{\Card}{\operatorname{Card}}
\providecommand{\E}{\operatorname{E}}
\providecommand{\V}{\operatorname{V}}
\providecommand{\degre}{^{\circ}}
\providecommand{\un}[1]{\,\mathrm{#1}}
\providecommand{\sci}[2]{#1\times10^{#2}}
\providecommand{\pgcd}{\operatorname{PGCD}}
\providecommand{\ppcm}{\operatorname{PPCM}}
\title{Dérivation, convexité et continuité}
\subtitle{Cours et exemples}
\date{}
\author{}
\begin{document}
\begin{frame}[plain]
\vfill
{\usebeamerfont{title}\inserttitle\par}
\medskip
{\usebeamerfont{subtitle}\insertsubtitle\par}
\medskip
{\usebeamercolor[fg]{progress bar}\rule{\textwidth}{.5pt}}
\vfill
\end{frame}
\begin{frame}{Objectifs}
Dériver une composée ; étudier variations et convexité ; démontrer existence et unicité ; approcher une solution.
\end{frame}
\begin{frame}{Dériver une composée}
Si \(u\) est dérivable sur \(I\) et \(v\) dérivable sur un intervalle contenant \(u(I)\), alors \((v\circ u)'=(v'\circ u)u'\). Les domaines doivent être vérifiés avant le calcul.
\[ (u^n)'=nu'u^{n-1},\quad(e^u)'=u'e^u,\quad
(\dfrac{1}{u})'=-\dfrac{u'}{u^2},\quad(\sqrt u)'=\dfrac{u'}{2\sqrt u}\quad(u>0).\]
Par exemple, \(f(x)=\sqrt{3x+1}\) est dérivable pour \(x>-\dfrac{1}{3}\) et \(f'(x)=\dfrac{3}{2\sqrt{3x+1}}\). Le signe de \(f'\) détermine les variations de \(f\) ; le signe de \(f\) ne les détermine pas.
\end{frame}
\begin{frame}{Convexité}
Une fonction \(f\) est convexe sur un intervalle \(I\) si sa courbe est sous ses cordes : pour \(x,y\in I\) et \(t\in[0,1]\),
\[f((1-t)x+ty)\le(1-t)f(x)+tf(y).\]
Pour une fonction deux fois dérivable, on admet les équivalences : \(f\) convexe ; \(f'\) croissante ; \(f''\ge0\) ; la courbe est au-dessus de toutes ses tangentes. Les inégalités sont inversées pour une fonction concave. Un point d'inflexion est un point où la courbe change de convexité.
\end{frame}
\begin{frame}{Convexité}
Supposons \(f''\ge0\). Fixons \(a\) et posons \(g(x)=f(x)-f(a)-f'(a)(x-a)\). Comme \(f'\) est croissante, \(g'(x)=f'(x)-f'(a)\le0\) pour \(x\le a\) et \(\ge0\) pour \(x\ge a\). Ainsi \(g\) atteint son minimum \(g(a)=0\) en \(a\), donc \(f(x)\ge f(a)+f'(a)(x-a)\).
\end{frame}
\begin{frame}{Convexité}
\(x^4\) est convexe et \(f''(0)=0\), mais \(0\) n'est pas une abscisse d'inflexion. \(x^3\) change de convexité en \(0\). La convexité de l'exponentielle et sa tangente en \(0\) donnent \(e^x\ge1+x\).
\end{frame}
\begin{frame}{Continuité et équations}
\(f\) est continue en \(a\) si \(\lim_{x\to a}f(x)=f(a)\). Toute fonction dérivable est continue. La réciproque est fausse : \(|x|\) est continue en \(0\) mais non dérivable en \(0\).
\end{frame}
\begin{frame}{Continuité et équations}
Si \(f\) est continue sur \([a,b]\), tout réel compris entre \(f(a)\) et \(f(b)\) possède au moins un antécédent dans \([a,b]\). Si \(f\) est en outre strictement monotone, cet antécédent est unique. On peut aussi travailler sur un intervalle ouvert en utilisant les limites aux bornes.
\end{frame}
\begin{frame}{Continuité et équations}
\(f(x)=x^3+x-1\) est continue sur \([0,1]\). Comme \(f(0)=-1\) et \(f(1)=1\), le TVI donne un zéro dans \(]0,1[\). Puisque \(f'(x)=3x^2+1>0\), il est unique. \(f(0{,}68)<0<f(0{,}69)\) fournit un encadrement, sans remplacer les justifications d'existence et d'unicité.
\end{frame}
\begin{frame}{Continuité et équations}
Pour \(x\ne1\), poser \(g(x)=\dfrac{x^2-1}{x-1}\), et poser \(g(1)=a\). Comme \(g(x)=x+1\) lorsque \(x\ne1\), \(\lim_{x\to1}g(x)=2\). La fonction est continue en \(1\) si et seulement si \(a=2\). La valeur au point doit donc être comparée à la limite, et non remplacée par celle-ci sans justification.
\end{frame}
\begin{frame}{Continuité et équations}
Si \(u_n\to\ell\) et \(f\) est continue en \(\ell\), alors \(f(u_n)\to f(\ell)\). Pour une suite récurrente convergente \(u_{n+1}=f(u_n)\), on obtient \(\ell=f(\ell)\).
\end{frame}
\begin{frame}{Dichotomie}
\begin{columns}[T,onlytextwidth]
\begin{column}{.57\textwidth}
Pour \(f\) continue et \(f(a)f(b)<0\), calculer \(m=\dfrac{a+b}{2}\). Si \(f(m)=0\), arrêter ; sinon garder la moitié dont les extrémités ont des images de signes contraires. Après \(n\) étapes, la longueur vaut \(\dfrac{b-a}{2^n}\). Le milieu donne une erreur au plus la demi-longueur. La méthode ne garantit pas l'unicité ; celle-ci se prouve séparément.   \(y=x^3-3x\) et sa tangente à l'inflexion.
\end{column}
\begin{column}{.4\textwidth}
\centering
\includegraphics[width=\linewidth,height=4.5cm,keepaspectratio]{figures/terminale-06-1.pdf}
\end{column}
\end{columns}
\end{frame}
\begin{frame}{Activité de recherche}
\textbf{Dérivées et lectures graphiques - 35 min.}

On considère la fonction \(f\) définie sur \(\R\) par \(f(x)=x^3-3x\). On note \(\mathcal C_f\) sa courbe représentative.
\end{frame}
\begin{frame}{Activité de recherche (2)}
\begin{enumerate}
\setcounter{enumi}{0}\item Calculer \(f' (x)\) et \(f'' (x)\), puis dresser leurs tableaux de signes.

\item En déduire les variations de \(f\) et ses intervalles de convexité et de concavité.

\item Déterminer les coordonnées du point d'inflexion de \(\mathcal C_f\) et une équation de la tangente en ce point.

\item En utilisant les limites, les extrema et le théorème des valeurs intermédiaires, déterminer le nombre de solutions réelles de \(f(x)=1\).

\item Expliquer, à partir de cet exemple, pourquoi la continuité ne suffit pas à assurer l'unicité d'une solution.

\end{enumerate}
\end{frame}
\end{document}
